% sections/13_verification.tex
\section{Verification, Test Coverage, and Formal Assurance}%
\label{sec:verification}

This section specifies the test suite that a conforming implementation of
the \ISI v\,1.0 pipeline must pass.  The suite is organised into five
categories: unit tests, property-based tests, golden-file tests,
integration tests, and integrity-chain tests.  Each test is designated
\textbf{T-\textit{n}} and mapped to the determinism invariant(s)
(\cref{sec:det:invariants}) it verifies.

\subsection{Test Infrastructure}\label{sec:test:infra}

The test suite is executed via a standard test runner (e.g., \texttt{pytest}).
All tests are deterministic: no test uses random seeds, network access, or
system clock queries.  Tests that exercise stochastic properties
(e.g., permutation invariance) use a fixed, reproducible set of permutations
drawn from a seeded pseudo-random generator whose seed is committed to the
test codebase.

Test data consists of two reference datasets:

\begin{description}[style=nextline]
  \item[\texttt{fixtures/reference\_small.json}]
    A synthetic dataset of five countries and six axes, with known
    closed-form solutions.  Used for unit tests and property-based tests
    where exact expected values can be computed by hand.
  \item[\texttt{fixtures/reference\_eu27.json}]
    The full EU-27 input dataset for vintage~2024.  Used for golden-file
    tests and integration tests.  The expected outputs are stored as golden
    files alongside the fixture.
\end{description}

\subsection{Unit Tests}\label{sec:test:unit}

Unit tests verify individual functions in isolation.  Each test provides
a known input and asserts exact equality with the expected output.

\begin{longtable}{@{}l >{\raggedright\arraybackslash}p{6cm} l@{}}
  \caption{Unit test registry.}\label{tab:test-unit} \\
  \toprule
  \textbf{ID} & \textbf{Description} & \textbf{Invariant} \\
  \midrule
  \endfirsthead
  \caption[]{Unit test registry (continued).} \\
  \toprule
  \textbf{ID} & \textbf{Description} & \textbf{Invariant} \\
  \midrule
  \endhead
  \midrule
  \multicolumn{3}{r@{}}{\footnotesize\itshape Continued on next page} \\
  \endfoot
  \bottomrule
  \endlastfoot
  \textbf{T-01} & \HHI of a uniform share vector
    \((\tfrac{1}{n}, \ldots, \tfrac{1}{n})\) equals \(\tfrac{1}{n}\). &
    --- \\
  \textbf{T-02} & \HHI of a degenerate share vector \((1, 0, \ldots, 0)\)
    equals~\(1\). & --- \\
  \textbf{T-03} & Dual-channel arithmetic mean: \(0.5 C^{(A)} + 0.5 C^{(B)}\)
    matches expected value for three test cases. & --- \\
  \textbf{T-04} & Category-weighted Channel~B (\cref{eq:weighted-b}) matches
    hand-computed value for two sub-categories with known volumes. & --- \\
  \textbf{T-05} & Fuel average (\cref{eq:fuel-avg}) matches expected value
    when one, two, and three fuel categories are present. & --- \\
  \textbf{T-06} & Composite (\cref{eq:composite}) matches arithmetic mean of
    six known axis scores. & --- \\
  \textbf{T-07} & Missing-axis composite (\cref{eq:composite-partial}) with
    \(|\mathcal{A}_i| = 4\) matches expected value. & \textbf{D-10} \\
  \textbf{T-08} & Rounding: input \texttt{"0.123456785"} at precision~8 returns
    \texttt{"0.12345678"} (ties to even). & \textbf{D-5} \\
  \textbf{T-09} & Rounding: input \texttt{"0.123456795"} at precision~8 returns
    \texttt{"0.12345680"} (ties to even). & \textbf{D-5} \\
  \textbf{T-10} & Classification of boundary values: \(0.50\) maps to
    \texttt{highly\_concentrated}; \(0.49999999\) maps to
    \texttt{moderately\_concentrated}. & \textbf{D-7} \\
  \textbf{T-11} & Classification of all four tier boundaries
    (\(0.15\), \(0.25\), \(0.50\), and values just below each). &
    \textbf{D-7} \\
  \textbf{T-12} & Scenario clamping: input axis score \(0.95\) with
    \(\alpha = +0.20\) returns \(1.0\), not \(1.14\). & \textbf{D-9} \\
  \textbf{T-13} & Scenario clamping: input axis score \(0.03\) with
    \(\alpha = -0.20\) returns \(0.024\), not a negative value. &
    \textbf{D-9} \\
  \textbf{T-14} & Ranking tie-break: two countries with identical composite
    scores are ordered by ISO alpha-2 code (ascending). & \textbf{D-8} \\
  \textbf{T-15} & Canonical hash string for a known country record matches
    a pre-computed SHA-256 digest. & \textbf{D-6} \\
\end{longtable}

\subsection{Property-Based Tests}\label{sec:test:property}

Property-based tests verify algebraic properties that must hold for
\emph{any} valid input, not merely for specific test vectors.  Inputs are
generated from constrained domains using a combinatorial or exhaustive
strategy over small state spaces.

\begin{longtable}{@{}l >{\raggedright\arraybackslash}p{5.5cm} l@{}}
  \caption{Property-based test registry.}\label{tab:test-property} \\
  \toprule
  \textbf{ID} & \textbf{Property} & \textbf{Invariant} \\
  \midrule
  \endfirsthead
  \caption[]{Property-based test registry (continued).} \\
  \toprule
  \textbf{ID} & \textbf{Property} & \textbf{Invariant} \\
  \midrule
  \endhead
  \midrule
  \multicolumn{3}{r@{}}{\footnotesize\itshape Continued on next page} \\
  \endfoot
  \bottomrule
  \endlastfoot
  \textbf{T-20} & \textbf{Unit-interval bound.}\enspace  For any valid share
    vector, \(0 < C_i^{(\mathrm{ch})} \leq 1\). & --- \\
  \textbf{T-21} & \textbf{Share-vector normalisation.}\enspace  Input share
    vectors that do not sum to~1 (within \(\epsilon = 10^{-9}\)) are rejected
    by the validation layer. & --- \\
  \textbf{T-22} & \textbf{Monotonicity of concentration.}\enspace  Shifting
    share mass from a small partner to the largest partner weakly increases
    the \HHI. & --- \\
  \textbf{T-23} & \textbf{Composite bounds.}\enspace
    \(0 < \mathrm{ISI}_i \leq 1\) for all countries with at least one
    computable axis. & --- \\
  \textbf{T-24} & \textbf{Classification monotonicity.}\enspace  If
    \(s_1 > s_2\), then
    \(\mathrm{tier}(s_1) \geq \mathrm{tier}(s_2)\) under the ordering
    \texttt{highly} \(>\) \texttt{moderately} \(>\) \texttt{mildly} \(>\)
    \texttt{unconcentrated}. & \textbf{D-7} \\
  \textbf{T-25} & \textbf{Ranking antisymmetry.}\enspace  If country~\(i\)
    ranks above country~\(j\), then country~\(j\) does not rank above
    country~\(i\). & \textbf{D-8} \\
  \textbf{T-26} & \textbf{Ranking totality.}\enspace  Every country in the
    snapshot receives exactly one rank.  No rank is assigned to two countries
    unless their composite scores are identical (in which case the tie-break
    applies). & \textbf{D-8} \\
  \textbf{T-27} & \textbf{Scenario bounds.}\enspace  For any valid baseline
    score and any \(\alpha \in [-0.20, +0.20]\), the simulated score lies
    in \([0, 1]\). & \textbf{D-9} \\
  \textbf{T-28} & \textbf{Scenario identity.}\enspace  When
    \(\alpha_a = 0\) for all~\(a\), the simulated scores equal the baseline
    scores exactly. & \textbf{D-9} \\
  \textbf{T-29} & \textbf{Rounding idempotency.}\enspace
    \(R_8(R_8(x)) = R_8(x)\) for all \(x\) in the domain. & \textbf{D-5} \\
  \textbf{T-30} & \textbf{Hash stability.}\enspace  The SHA-256 digest of
    a canonical string is invariant across invocations. & \textbf{D-6} \\
\end{longtable}

\subsection{Golden-File Tests}\label{sec:test:golden}

Golden-file tests compare the full pipeline output against a pre-computed
reference.  They are the primary mechanism for verifying end-to-end
determinism.

\begin{longtable}{@{}l >{\raggedright\arraybackslash}p{5.5cm} l@{}}
  \caption{Golden-file test registry.}\label{tab:test-golden} \\
  \toprule
  \textbf{ID} & \textbf{Description} & \textbf{Invariant} \\
  \midrule
  \endfirsthead
  \caption[]{Golden-file test registry (continued).} \\
  \toprule
  \textbf{ID} & \textbf{Description} & \textbf{Invariant} \\
  \midrule
  \endhead
  \midrule
  \multicolumn{3}{r@{}}{\footnotesize\itshape Continued on next page} \\
  \endfoot
  \bottomrule
  \endlastfoot
  \textbf{T-40} & \textbf{Small-dataset golden file.}\enspace  Pipeline
    output on \texttt{reference\_small.json} matches the golden file
    byte-for-byte. & \textbf{D-1}, \textbf{D-2} \\
  \textbf{T-41} & \textbf{EU-27 golden file.}\enspace  Pipeline output on
    \texttt{reference\_eu27.json} matches the golden file byte-for-byte.
    This test verifies the production dataset. & \textbf{D-1}, \textbf{D-2} \\
  \textbf{T-42} & \textbf{Idempotency.}\enspace  Running the pipeline twice
    on \texttt{reference\_eu27.json} produces byte-identical output. &
    \textbf{D-1} \\
  \textbf{T-43} & \textbf{Row-permutation invariance.}\enspace  Running the
    pipeline on 10~distinct random permutations of the input rows of
    \texttt{reference\_small.json} produces identical output each time. &
    \textbf{D-4} \\
  \textbf{T-44} & \textbf{Scenario golden file.}\enspace  Running the
    scenario engine with a fixed adjustment vector on a fixed country
    produces output matching a pre-computed golden file. &
    \textbf{D-9} \\
\end{longtable}

\subsection{Integration Tests}\label{sec:test:integration}

Integration tests verify the end-to-end pipeline including data ingestion,
validation, computation, serialisation, and integrity-chain construction.

\begin{longtable}{@{}l >{\raggedright\arraybackslash}p{5.5cm} l@{}}
  \caption{Integration test registry.}\label{tab:test-integration} \\
  \toprule
  \textbf{ID} & \textbf{Description} & \textbf{Invariant} \\
  \midrule
  \endfirsthead
  \caption[]{Integration test registry (continued).} \\
  \toprule
  \textbf{ID} & \textbf{Description} & \textbf{Invariant} \\
  \midrule
  \endhead
  \midrule
  \multicolumn{3}{r@{}}{\footnotesize\itshape Continued on next page} \\
  \endfoot
  \bottomrule
  \endlastfoot
  \textbf{T-50} & \textbf{Snapshot completeness.}\enspace  The materialised
    snapshot directory contains all mandatory files listed in
    \cref{sec:snapshot:directory}. & --- \\
  \textbf{T-51} & \textbf{Manifest consistency.}\enspace  Every file listed
    in \texttt{MANIFEST.json} exists in the directory, and its SHA-256 hash
    matches the recorded value. & \textbf{D-6} \\
  \textbf{T-52} & \textbf{Hash-summary consistency.}\enspace  Recomputing
    all country hashes from \texttt{country\_scores.csv} yields the same
    values as \texttt{HASH\_SUMMARY.json}. & \textbf{D-6} \\
  \textbf{T-53} & \textbf{Snapshot-hash consistency.}\enspace  The snapshot
    hash \(H_t\) recomputed from sorted country hashes matches the value in
    \texttt{HASH\_SUMMARY.json}. & \textbf{D-6} \\
  \textbf{T-54} & \textbf{Signature verification.}\enspace  The Ed25519
    signature in \texttt{SIGNATURE.json} verifies against \(H_t\) using the
    published public key. & --- \\
  \textbf{T-55} & \textbf{CSV schema compliance.}\enspace  All CSV output
    files conform to the column schemas specified in
    \cref{sec:snapshot:output}. & --- \\
  \textbf{T-56} & \textbf{JSON schema compliance.}\enspace  The methodology
    registry and snapshot metadata files contain all mandatory keys listed in
    \cref{sec:snapshot:registry,sec:snapshot:metadata}. & --- \\
  \textbf{T-57} & \textbf{Score-string format.}\enspace  All score strings
    in CSV and JSON output files match the regex
    \texttt{\textasciicircum{}(null|-?[0-9]+\textbackslash.[0-9]\{8\})\$}. & \textbf{D-5} \\
  \textbf{T-58} & \textbf{Warning propagation.}\enspace  Warning codes
    generated during computation appear in \texttt{warnings.csv} and are
    correctly associated with the respective country--axis pair. & --- \\
  \textbf{T-59} & \textbf{Fallback-basis recording.}\enspace  The
    \texttt{fallback\_basis.csv} file records one entry per country--axis
    pair, and the recorded basis matches the channel availability in the
    input data. & --- \\
\end{longtable}

\subsection{Integrity-Chain Tests}\label{sec:test:integrity}

These tests specifically verify the cryptographic integrity chain in
isolation.

\begin{longtable}{@{}l >{\raggedright\arraybackslash}p{5.5cm} l@{}}
  \caption{Integrity-chain test registry.}\label{tab:test-integrity} \\
  \toprule
  \textbf{ID} & \textbf{Description} & \textbf{Invariant} \\
  \midrule
  \endfirsthead
  \caption[]{Integrity-chain test registry (continued).} \\
  \toprule
  \textbf{ID} & \textbf{Description} & \textbf{Invariant} \\
  \midrule
  \endhead
  \midrule
  \multicolumn{3}{r@{}}{\footnotesize\itshape Continued on next page} \\
  \endfoot
  \bottomrule
  \endlastfoot
  \textbf{T-60} & \textbf{Country-hash determinism.}\enspace  The same
    country record always produces the same SHA-256 digest. &
    \textbf{D-6} \\
  \textbf{T-61} & \textbf{Country-hash sensitivity.}\enspace  Changing a
    single decimal digit in one axis score produces a different country
    hash. & \textbf{D-6} \\
  \textbf{T-62} & \textbf{Snapshot-hash order dependence.}\enspace  The
    snapshot hash changes if the country-hash concatenation order is
    altered (verifying that lexicographic sorting is enforced). &
    \textbf{D-6} \\
  \textbf{T-63} & \textbf{Signature validity.}\enspace  A freshly computed
    signature over \(H_t\) verifies with the corresponding public key. &
    --- \\
  \textbf{T-64} & \textbf{Tamper detection.}\enspace  Modifying one byte in
    any score file causes the manifest hash to differ from the recorded
    value, triggering a verification failure. & \textbf{D-6} \\
  \textbf{T-65} & \textbf{Signature rejection.}\enspace  An Ed25519
    signature produced with a different private key fails verification
    against the published public key. & --- \\
\end{longtable}

\subsection{Coverage Requirements}\label{sec:test:coverage}

A conforming implementation must satisfy the following minimum coverage
thresholds:

\begin{table}[htbp]
  \centering\small
  \caption{Minimum test-coverage requirements.}\label{tab:test-coverage}
  \begin{tabular}{@{}l l l@{}}
    \toprule
    \textbf{Module} & \textbf{Line coverage} & \textbf{Branch coverage} \\
    \midrule
    Concentration scoring        & 100\% & 100\% \\
    Cross-channel aggregation    & 100\% & 100\% \\
    Composite averaging          & 100\% & 100\% \\
    Classification               & 100\% & 100\% \\
    Rounding                     & 100\% & 100\% \\
    Ranking                      & 100\% & 100\% \\
    Scenario engine              & 100\% & 100\% \\
    Integrity chain (hashing)    & 100\% & 100\% \\
    Integrity chain (signing)    & 100\% & 100\% \\
    Serialisation (JSON/CSV)     & \(\geq 95\%\) & \(\geq 90\%\) \\
    Data ingestion / parsing     & \(\geq 90\%\) & \(\geq 85\%\) \\
    Validation                   & \(\geq 95\%\) & \(\geq 90\%\) \\
    \bottomrule
  \end{tabular}
\end{table}

All core mathematical and classification modules require 100\% line and
branch coverage.  This is non-negotiable: any untested branch in a scoring
or classification function represents a potential determinism violation.

\subsection{Regression Protocol}\label{sec:test:regression}

When a bug is discovered in the pipeline, the following protocol applies:

\begin{enumerate}
  \item \textbf{Bug isolation.}  A minimal reproducing test case is
    constructed and added to the test suite as a new test
    (\textbf{T-\textit{n}}).
  \item \textbf{Fix implementation.}  The bug is fixed in the pipeline
    codebase.
  \item \textbf{Golden-file update.}  If the fix changes the output for the
    reference datasets, the golden files are regenerated and the new expected
    output is committed alongside the fix.
  \item \textbf{Version increment.}  If the fix changes published output
    (axis scores, composite scores, classifications, or hashes), a minor or
    major version increment is triggered under the governance protocol
    (\cref{sec:governance}).
  \item \textbf{Non-regression gate.}  The full test suite must pass before
    the fix is merged.  No test may be deleted to accommodate a fix.
\end{enumerate}

\subsection{Continuous Integration}\label{sec:test:ci}

The test suite runs automatically on every commit to the main branch and on
every pull request.  The CI pipeline enforces:

\begin{enumerate}
  \item All tests pass (zero failures, zero errors).
  \item Coverage thresholds (\cref{tab:test-coverage}) are met.
  \item No new static-analysis warnings are introduced (linting via
    \texttt{ruff} or equivalent).
  \item No new type-checking errors are introduced (type checking via
    \texttt{mypy} in strict mode or equivalent).
\end{enumerate}

A pull request that fails any of these gates cannot be merged.

\subsection{Formal Assurance Statement}\label{sec:test:assurance}

The test suite does not constitute a formal proof of correctness.  It
provides a high-confidence empirical assurance that the pipeline conforms
to its specification.  The assurance is bounded by the coverage and
diversity of the test inputs.  Specifically:

\begin{itemize}
  \item Property-based tests (\textbf{T-20} through \textbf{T-30}) verify
    algebraic invariants over parameterised input spaces, providing stronger
    guarantees than fixed test vectors alone.
  \item Golden-file tests (\textbf{T-40} through \textbf{T-44}) provide
    byte-level regression protection for the complete pipeline.
  \item Integrity-chain tests (\textbf{T-60} through \textbf{T-65}) verify
    the cryptographic binding between scores and published hashes.
\end{itemize}

Formal verification (e.g., proof-carrying code, model checking) is not
applied in v\,1.0 but is identified as a candidate for future methodology
versions where the state space and cost of verification warrant it.
